\section{Discussion}
\label{sec:discussion}

% ZNTODO: write text for these points flagged elsewhere in the manuscript
% (where they are flagged is given in brackets; each place has a `[TODO ... forward ref]` or `ZNTODO` to update once the text is written):
% - more robust interpolation methods
%   (methods, general approach, "More robust interpolation methods are an area for future work")
% - harmonisation of ice core data to the observation network is crude
%   (methods, N_2O and CO_2 Step 4, "This is crude and a clear area for further investigation")
% - the configuration error which means we don't use the Menking et al. (2025) CH_4 data,
%   and the resulting care needed when using the dataset for pre-1850 studies
%   (methods, CH_4 Step 4; output requirements, "we didn't use Menking dataset for pre-1850")
% - calibration scales are not converted and observations on different scales are combined
%   (methods, CH_4; output requirements, "we don't consider different scales of observations").
%   Partly covered by the fifth point below.
% - constant poleward extrapolation as a clear weakness
%   (methods, CFC12-like Step 2, "[TODO note this as a clear weakness in discussion]")
% - needing more rapidly updated historical emissions estimates, or an approach without the regression against emissions;
%   the need for a method which creates a complete dataset back to year 1
%   and can be projected forward from emissions and MAGICC alone (all we have when creating scenario forcings)
%   (methods, CFC12-like Step 4)
% - fixing the upstream processing so it doesn't produce small negative values.
%   Probably covered by the sixth point below (scaling down the latitudinal gradient and seasonality).
% - the CH_4 seasonal cycle may be slightly too strong; a more sophisticated handling than relative seasonality
%   (results, CH_4, "[TODO add section on this in discussion]")
% - C_8F_18 should be considered for future iterations
%   (methods, C_8F_18, "should be considered for future iterations")
% - annual updates of the forcings (abstract, "discuss annual updates more clearly in discussion")
%   and a REF for forcings (output requirements, "put points below into discussion").
%   Annual updates ameliorate the risk of not delivering on time;
%   a REF would help validation against the experiment design and speed up finding bugs.
%   NOTES.md's "How to do extensions" sketches the approach:
%   optimise the existing latitudinal gradient/seasonality change against new network data
%   rather than re-deriving or regressing PCs.
% - vertical resolution (see the `ZNTODO` at the end of this section; methods, Vertical dimension)
% - gathering more data sources, e.g. the Nevado Huascar\'an ice core (see the `ZNTODO` at the end of this section)
% - piControl should be done by repeating the 1850 values: make sure this is stated somewhere
%   (output requirements `ZNTODO`, may belong in the results rather than here)

The output dataset shows differences from its inputs, other observation networks and the CMIP6 dataset
that are small in the context of earth system modelling (Section~\ref{sec:results}).
% value-check: {"tag": "co2-diff-from-cmip6-all-radiative-effect", "unit": "W/m^2", "upper": 0.06, "lower": 0}
% value-check: {"tag": "ch4-diff-from-cmip6-all-radiative-effect", "unit": "W/m^2", "upper": 0.06, "lower": 0}
% value-check: {"tag": "n2o-diff-from-cmip6-all-radiative-effect", "unit": "W/m^2", "upper": 0.06, "lower": 0}
% value-check: {"tag": "cfc12-diff-from-cmip6-all-radiative-effect", "unit": "W/m^2", "upper": 0.06, "lower": 0}
% value-check: {"tag": "cfc12eq-diff-from-cmip6-all-radiative-effect", "unit": "W/m^2", "upper": 0.06, "lower": 0}
% value-check: {"tag": "hfc134aeq-diff-from-cmip6-all-radiative-effect", "unit": "W/m^2", "upper": 0.06, "lower": 0}
For example, in approximate radiative effect terms, the differences from the CMIP6 dataset are less than 0.06~W/m2 over the length of the dataset
% value-check: {"tag": "co2-diff-from-cmip6-1850-on-radiative-effect", "unit": "W/m^2", "upper": 0.03, "lower": 0}
% value-check: {"tag": "ch4-diff-from-cmip6-1850-on-radiative-effect", "unit": "W/m^2", "upper": 0.03, "lower": 0}
% value-check: {"tag": "n2o-diff-from-cmip6-1850-on-radiative-effect", "unit": "W/m^2", "upper": 0.03, "lower": 0}
% value-check: {"tag": "cfc12-diff-from-cmip6-1850-on-radiative-effect", "unit": "W/m^2", "upper": 0.03, "lower": 0}
% value-check: {"tag": "cfc12eq-diff-from-cmip6-1850-on-radiative-effect", "unit": "W/m^2", "upper": 0.03, "lower": 0}
% value-check: {"tag": "hfc134aeq-diff-from-cmip6-1850-on-radiative-effect", "unit": "W/m^2", "upper": 0.03, "lower": 0}
and less than 0.03~W/m2 over the period used for CMIP modelling (1850 onwards).
There are nonetheless a number of areas for potential improvement.

The first is a lack of new input datasets compared to CMIP6.
This is particularly the case for a number of gases with small radiative forcings.
We suggest picking a smaller, more clearly supported set of gases
or aligning more closely with the gases included in an exercise like IGCC,
which has committed to annual updates.

The second is that we do not use satellite data as an input dataset.
While satellite information is only available for a limited number of gases and a limited time period,
it provides a completely different set of spatiotemporal characteristics than the ground-based network.
The inclusion of information from satellite data could therefore enhance the available input dataset
and improve the output data product,
particularly by reducing the extent to which we have to rely on spatial interpolation that comes with high uncertainty.

The third is that we do not report uncertainties in the output dataset.
While we do not expect these uncertainties to be large in the context of earth system modelling,
they would help contextualise the dataset in a quantified way that is currently lacking.
For example, the inclusion of uncertainties would likely make clear that the leading uncertainty
is not related to measurement error.
Instead, it is methodological: the current method requires a number of modelling steps/assumptions
to go from the sparse observation networks we have available to a spatially and temporally complete dataset
over the entire earth's surface from year 1 to approximately present-day.
This error will also vary strongly over time:
far less modelling is required to make a spatially complete estimate in the period where we have observation networks
than when our only estimates come from one or two ice core records.
% IGCC's values: Sect. 3 of Forster et al. (2026), which carries over the AR6 uncertainties (checked 2026-10-09).
% Our CH_4 difference from NOAA: checked in the results (tag `ch4-monthly-diff-from-noaa`).
It is also not clear that the uncertainties which are reported elsewhere capture this methodological uncertainty.
For example, IGCC \citep{forster_igcc_2026} reports uncertainties in global-mean surface concentrations
of 0.4~ppm for CO_2, 3.3~ppb for CH_4 and 0.4~ppb for N_2O.
These appear too small, particularly for CH_4 and N_2O,
given that the underlying observation networks do not have global coverage.
% ZNTODO: copy value check from results in here
Our CH_4 global-mean differs from NOAA's by around 10~ppb (Section~\ref{ssec:results-ch4}),
despite both being based on largely the same observations.
In our experience, this understanding is generally missing
and quantified uncertainties would help a) inform users of the origin of uncertainty,
b) encourage observation networks to engage more directly with the question of how to reduce the error
that is introduced between raw observations and an estimate of spatially and temporally complete values
and c) identify where effort should be directed to reduce uncertainty in greenhouse gas forcing datasets.

The fourth is that we, rather than the observation networks, are producing the forcings.
As we understand it, to date, the observation networks and earth system modelling communities have had relatively little direct contact.
Instead, teams that produce forcings (like us) have acted as `bridges' between the two:
translating observation products into the spatiotemporally complete datasets required for use as model inputs (i.e. forcings).
While we have developed expertise in acting as this bridge,
we are not experts in either the raw observations or earth system modelling.
There is scope for investigating whether better outcomes would be achieved if the observation and earth system modelling communities worked more closely.
As a readily available example, the AGAGE community already produces a model-based reanalysis
that provides global coverage with spatial and temporal variation.
This could make our analysis redundant and the earth system modelling community could use this AGAGE product (or a variation thereof) directly.
To date, no investigation has been done into the pros and cons of such a switch
and the resources to do such an investigation have not been available,
but our suspicion is that we would get a better result if the AGAGE team produced a product that matched ESM needs directly,
rather than having us act as an intermediary.
Such an investigation would have to consider a number of questions,
for example whether modelling-based or statistical approaches to infilling are better for the intended use cases
(given that modelling-based approaches are generally more grounded in physics, but are also more prone to producing extreme results),
how observation providers could expand their data products to provide the required temporal coverage
and which observation networks/data sources should be used given that there can be more than one provider and forcings are usually derived based on a composite of sources.

The fifth is that we do not convert between the calibration scales of the different observation networks,
nor have we explored the differences between them in depth.
For the gases with the largest radiative forcing, the effect is negligible:
we only use data from NOAA for CO_2,
the difference between the AGAGE and NOAA scales for CH_4 is less than 0.1\%
(\citealp{meinshausen_historical-ghgs_2017} converted AGAGE CH_4 to the NOAA scale by multiplying by 1.0003)
and the AGAGE and NOAA scales for N_2O are compatible without conversion
\citep{meinshausen_historical-ghgs_2017}.
For halocarbons, the differences between calibration scales are small
and these gases have smaller radiative forcings.
A more careful treatment of calibration scales,
ideally developed together with the observation networks,
would be an improvement.

The sixth is internal consistency issues within our method, which lead to the use of scaling to avoid unphysical negative values.
This is most visible for the CFC12-like and C_4F_10-like gases,
where we have to scale down the latitudinal gradient and the seasonality to avoid negative values in the output
(Sections~\ref{sssec:methods-cfc12-like-prepare-components} and \ref{sssec:methods-c4f10-like-prepare-components}).
The need for this points to three inconsistencies.
The first is between our pre-industrial values and the emissions used to extend the latitudinal gradient.
The global-, annual-mean is held at its pre-industrial value until the pre-industrial year and only rises slowly after it,
whereas the latitudinal gradient follows emissions, which can already be well above zero in these years.
A better approach would derive the components in a way that guarantees their consistency,
for example by constraining the latitudinal gradient and the seasonality relative to the global-mean when they are derived and extended
or by considering the consistency between the emissions estimates and concentration changes.
% % zalue-check: {"tag": "cfc12-like-lat-gradient-capped-years-after-pre-industrial", "unit": "yr", "lower": 44, "upper": 44}
% % zalue-check: {"tag": "c4f10-like-lat-gradient-capped-years-first", "unit": "yr", "lower": 1934, "upper": 1934}
% % zalue-check: {"tag": "c4f10-like-lat-gradient-capped-years-last", "unit": "yr", "lower": 2001, "upper": 2001}
% As a result, we scale down the latitudinal gradient of nearly every CFC12-like gas
% in the (up to 44) years after its pre-industrial year
% and of the C_4F_10-like gases in some months between 1934 and 2001.
The second is specific to HFC-152a,
% value-check: {"tag": "hfc152a-lat-gradient-capped-years-first", "unit": "yr", "lower": 1991, "upper": 1991}
% value-check: {"tag": "hfc152a-lat-gradient-capped-years-last", "unit": "yr", "lower": 2020, "upper": 2020}
for which we scale down the latitudinal gradient in some months between 1991 and 2020.
This period includes the years in which we have observations,
so the extension back in time cannot be the only cause:
the latitudinal gradient we derive from the observations is itself too strong for the global-mean in these months.
This suggests that we need to consider how we handle the input data more carefully,
for example removing outliers or pollution events (Figure~\ref{fig:methods-hfc152a}) before performing PCA.
The third is specific to HFC-236fa,
% value-check: {"tag": "hfc236fa-seasonality-capped-years-first", "unit": "yr", "lower": 1996, "upper": 1996}
% value-check: {"tag": "hfc236fa-seasonality-capped-years-last", "unit": "yr", "lower": 2022, "upper": 2022}
% value-check: {"tag": "max-seasonality-fraction", "unit": "dimensionless", "lower": 0.35, "upper": 0.35}
for which the seasonality we derive from the observations is more than 35\% of the global-, annual-mean
in every year from 1996 to 2022.
A seasonal cycle of this size is not plausible for such a long-lived gas
and most likely reflects noise in the measurements of its very low concentrations (Figure~\ref{fig:methods-hfc236fa}).
As for HFC-152a, a more careful handling of the input data would likely solve this issue.
% % zalue-check: {"tag": "hfc236fa-lat-gradient-first-eof-variance", "unit": "dimensionless", "lower": 0.75, "upper": 0.85}
% (HFC-236fa is also the only CFC12-like gas for which the first latitudinal gradient EOF explains less than 90\% of the variance,
% Section~\ref{sssec:methods-cfc12-like-pca}).
Scaling the components down is a crude fix which future work should look to improve on.

% ZNTODO: vertical resolution
% ZNTODO: note about gathering more data sources e.g. Nevado Huascar\'an ice core which suggests an equitorial maximum for CH4,
% very different to current data
